% -----------------------------------------------
% Composer Fingerprinting: A Multi-Layer Feature Approach for Lieder Authorship Attribution
% ISMIR 2026 Submission
% -----------------------------------------------

\documentclass{article}
\usepackage[T1]{fontenc}
\usepackage[utf8]{inputenc}
\usepackage[submission]{ismir}
\usepackage{amsmath,cite,url}
\usepackage{graphicx}
\usepackage{color}
\usepackage{booktabs}
\usepackage{multirow}

% Title
\title{Composer Fingerprinting: A Multi-Layer Feature Approach for Lieder Authorship Attribution}

% Anonymous for double-blind review
\oneauthor
  {Anonymous Authors}
  {Anonymous Affiliations\\\texttt{anonymous@ismir.net}}

% CC license author names (abbreviated)
\def\authorname{Y. Wu and J. Pouliou}

\sloppy

\begin{document}

\maketitle

\begin{abstract}
Classifying musical compositions by their composer remains a challenging task in computational musicology, particularly when distinguishing between composers working within the same historical period and genre. This exploratory study investigates whether individual composer style can be captured through computational analysis of symbolic music representations in German Lieder. We propose a three-layer feature framework combining tonal tension (computed via the Spiral Array Model), harmonic complexity (measured as pitch class entropy), and pianistic texture (operationalized as onset density) to classify 264 Lieder by Franz Schubert, Robert Schumann, and Johannes Brahms. Our experiments demonstrate that carefully designed handcrafted features (54 dimensions, excluding note count and velocity to avoid confounding variables) achieve approximately 67\% balanced accuracy using a Support Vector Machine classifier with feature selection (top 20 features). This significantly outperforms both SVM (47.1\%) and MLP (45.3\%) classifiers on 768-dimensional pretrained transformer embeddings from Adversarial-MidiBERT on this limited corpus. Feature importance analysis reveals that unison ratio, texture variation, and pitch standard deviation are the most discriminative attributes. ANOVA results confirm significant between-composer variance in texture and melody features ($p < 0.05$). These findings suggest that domain-specific, interpretable features encoding musicological knowledge may be more effective than general-purpose pretrained representations when working with small, specialized corpora.
\end{abstract}

\section{Introduction}\label{sec:introduction}

Classifying musical compositions by their composer has been a longstanding challenge in computational musicology. While significant progress has been made in distinguishing between different historical eras or genres, accurately identifying individual composers working within the same period and stylistic tradition remains difficult even with modern machine learning techniques \cite{Alvarez2024}. This challenge is particularly pronounced in the German Lieder tradition of the Romantic era, where composers such as Franz Schubert (1797--1828), Robert Schumann (1810--1856), and Johannes Brahms (1833--1897) shared a common musical vocabulary while developing distinct individual styles.

The fundamental question underlying this research is whether computational methods can capture and quantify the subtle stylistic markers that musicologists recognize intuitively---the composer's ``fingerprint'' that distinguishes Schubert's lyrical simplicity from Schumann's poetic intensity and Brahms' classical restraint. Early approaches to this problem relied on statistical characterization of musical surfaces \cite{Youngblood1958}, while more recent work has employed sophisticated machine learning architectures including Support Vector Machines, Hidden Markov Models, and deep neural networks. However, many of these studies focus on era-level classification or use features that lack direct musicological interpretability.

This project adopts an \textit{exploratory} perspective: rather than claiming definitive composer attribution, we investigate which computational features show promise for capturing stylistic individuality and how these features align with musicological understanding of each composer's approach to the Lieder genre.

\subsection{Research Questions}

This exploratory study is guided by three research questions:

\textbf{RQ1:} Can symbolic features capture composer-specific patterns in German Lieder?

\textbf{RQ2:} Which musical dimensions (harmony, texture, melody) are most discriminative for composer classification?

\textbf{RQ3:} How do handcrafted, domain-specific features compare to pretrained transformer representations for composer classification with limited training data?

\subsection{Contributions}

This paper makes the following contributions:

\begin{enumerate}
    \item \textbf{Multi-layer Feature Framework:} We propose a theoretically-grounded approach combining three complementary feature categories---tonal tension (Spiral Array Model), harmonic complexity (pitch class entropy), and pianistic texture (onset density)---with detailed musicological justification for each design choice.
    
    \item \textbf{Systematic Comparison:} We present a comprehensive comparison of handcrafted features (12D, 54D) against pretrained transformer embeddings (768D MidiBERT) with both SVM and MLP classifiers, with results showing handcrafted features achieve ~67\% balanced accuracy (with feature selection) versus 47.1\% (SVM) and 45.3\% (MLP) for MidiBERT on our 264-piece corpus.
    
    \item \textbf{Feature Importance Analysis:} Using Random Forest importance ranking and ANOVA, we identify the most discriminative features (unison ratio, texture variation, pitch standard deviation) and provide musicological interpretation of why these features capture composer-specific patterns.
    
    \item \textbf{Methodological Rigor:} We explicitly exclude note count (confounding variable) and velocity features (editorial bias) from our analysis, demonstrating that classification can achieve meaningful accuracy based on genuine musical features alone.
    
    \item \textbf{Reproducible Pipeline:} We release all code, extracted features, and experimental configurations to support reproducible research in computational composer classification.
\end{enumerate}

\section{Background and Related Work}\label{sec:background}

\subsection{The German Lieder Tradition}

The German Lied (art song) represents one of the most intimate and expressive genres of Western classical music. Emerging in the late 18th century and flourishing throughout the Romantic period (1800--1900), the Lied combines poetry and music in a partnership between voice and piano.

\textbf{Franz Schubert} stands as the foundational figure of the German Lied tradition, composing more than 600 Lieder. Musicologically, Schubert's Lieder are characterized by lyrical melody, guitar-like arpeggiated accompaniment, text-driven harmony, and flexible formal approaches \cite{Rosen1995}.

\textbf{Robert Schumann} represents the high Romantic approach to Lieder composition. His style is distinguished by piano independence, arpeggiated textures, psychological depth, and cyclic unity across song cycles \cite{Daverio1997}.

\textbf{Johannes Brahms} composed nearly 200 Lieder while maintaining a conscious connection to classical forms. His approach reflects conservative harmony, dense texture, folk song influence, and classical restraint \cite{Musgrave1985}.

The shared tradition of these three composers presents a particular challenge for computational classification, as all worked within the same tonal system, similar formal conventions, comparable instrumental forces, and related poetic traditions.

\subsection{Computational Composer Classification}

The statistical characterization of musical style dates to Youngblood's \cite{Youngblood1958} pioneering work. Simonetta \cite{Simonetta2025} traces the evolution from early statistical methods to modern computational approaches, noting the shift from purely statistical descriptions to machine learning-based classification in the 1990s.

Recent work has explored diverse feature representations: Tonal Interval Vector (TIV) frameworks \cite{Herremans2016}, melodic contour analysis \cite{Pollastri2003}, and texture modeling \cite{Giraud2014}. Machine learning architectures have evolved from Hidden Markov Models \cite{HMM2001} to Support Vector Machines \cite{Alvarez2024} and deep neural networks.

\subsection{Deep Learning vs. Handcrafted Features}

The success of transformer-based models inspired similar approaches for symbolic music, including MusicBERT \cite{Chen2021}, MidiBERT-Piano \cite{Tsay2021}, and Adversarial-MidiBERT \cite{Zhao2025}. However, classical music corpora present unique challenges: limited availability, high annotation cost, and class imbalance. This creates a scenario where handcrafted features encoding domain knowledge may outperform general pretrained representations.

Feature extraction libraries such as jSymbolic \cite{McKay2006} and musif \cite{Llorens2023} facilitate automated feature extraction, though our approach differs in its theoretical grounding: each feature is selected based on musicological hypotheses about composer style.

\section{Methodology}\label{sec:methodology}

\subsection{Dataset Construction}

Our dataset comprises 264 German Lieder drawn from the OpenScore Lieder Repository \cite{OpenScore2023} and the Schubert Winterreise Dataset \cite{Winterreise2022}. Table~\ref{tab:dataset} shows the corpus composition.

\begin{table}[h]
\centering
\caption{Dataset composition by composer.}
\label{tab:dataset}
\begin{tabular}{lrrr}
\toprule
Composer & Pieces & Percentage & Date Range \\
\midrule
Franz Schubert & 84 & 31.8\% & 1815--1828 \\
Johannes Brahms & 109 & 41.3\% & 1852--1896 \\
Robert Schumann & 71 & 26.9\% & 1840--1852 \\
\midrule
\textbf{Total} & \textbf{264} & \textbf{100\%} & \textbf{1815--1896} \\
\bottomrule
\end{tabular}
\end{table}

All scores underwent preprocessing: format conversion (MusicXML to MIDI), voice separation using music21 \cite{Cuthbert2010} and partitura, quality filtering, meter detection, and composer name normalization.

\subsection{Feature Design Rationale}

Our feature design is grounded in musicological theory about composer style in the German Lieder tradition.

\subsubsection{Tonal Tension}

The Spiral Array Model \cite{Herremans2016,Herremans2019} provides a geometric representation of tonal space. We hypothesize that each composer develops individual patterns of harmonic tension and resolution. For each metric grid position, we collect sounding pitches, map them to Spiral Array 3D coordinates, compute the centroid, and calculate the mean Euclidean distance as tonal tension. Features: \texttt{tt\_mean}, \texttt{tt\_std}, \texttt{tt\_entropy}.

\subsubsection{Harmonic Complexity}

Pitch class distribution entropy measures chromatic saturation \cite{Kaliakatsos2011}. We hypothesize that Brahms' conservative harmonic language yields lower entropy than Schumann's chromatic explorations. For each beat, we compute Shannon entropy over 12 pitch classes. Features: \texttt{hc\_mean}, \texttt{hc\_std}, \texttt{hc\_entropy}.

\subsubsection{Pianistic Texture}

Texture in Lieder accompaniment ranges from sparse to dense \cite{Giraud2014}. We hypothesize that Schubert's guitar-like arpeggiation yields lower simultaneity than Brahms' chordal writing. For each beat in the piano part, we count simultaneous note onsets. Features: \texttt{pt\_mean}, \texttt{pt\_std}, \texttt{pt\_entropy}.

\subsubsection{Melodic Contour}

Melodic shape carries significant stylistic information \cite{Pollastri2003}. We extract the vocal line as a sequence of MIDI pitches and compute interval statistics. Features: \texttt{mc\_mean}, \texttt{mc\_std}, \texttt{mc\_entropy}.

\subsection{Feature Extraction Pipeline}

Feature extraction aligns to the metric grid rather than fixed time windows (quarter-note for 2/4, 3/4, 4/4; eighth-note for 6/8). Each bar is represented as a matrix of feature values, then aggregated into piece-level statistics (mean, standard deviation, entropy), yielding a 12-dimensional feature vector per piece.

\subsection{Handmade Feature Extension (54D)}

To complement the theory-driven 12D feature set, we extracted 54 additional statistical features across four musical dimensions: pitch features (f2--f10), rhythm/duration features (f16--f23), interval features (f24--f30, f52--f60), texture features (f47--f50), and higher-order statistics (f31--f34, f35--f46).

\subsubsection{Excluded Features}

\textbf{Note Count (f1)}: Initially included but removed from final analysis. Note count serves as a proxy for piece length, which may confound stylistic analysis with formal/structural choices unrelated to composer-specific musical language. By excluding this feature, we test whether genuine musical features (interval patterns, texture, rhythm) can achieve competitive classification performance.

\textbf{Velocity Features (f11--f15)}: MIDI velocity values in symbolic datasets often reflect editorial conventions of score transcribers rather than composer intent, as historical scores from the Romantic era specify dynamics qualitatively (e.g., \textit{p}, \textit{f}) rather than as numerical values (1--127). To ensure our model captures genuine stylistic patterns rather than data source artifacts, we excluded velocity features from our primary analysis.

\subsection{MidiBERT Feature Extraction}

For comparison, we extracted 768-dimensional embeddings using Adversarial-MidiBERT \cite{Zhao2025}. Each note is encoded as an 8-tuple (TimeSig, Tempo, Bar, Position, Instrument, Pitch, Duration, Velocity), processed through the pretrained transformer, and average-pooled to produce a single 768-dimensional vector per piece.

\subsection{Classification Framework}

We employ the following classifiers:

\textbf{Support Vector Machines (SVM)} with RBF kernel ($C=5.0$, \texttt{class\_weight='balanced'}). Evaluation uses 5-fold stratified cross-validation with balanced accuracy as the primary metric:

\begin{equation}\label{eq:balanced_acc}
\text{Balanced Accuracy} = \frac{1}{K} \sum_{k=1}^{K} \frac{TP_k}{TP_k + FN_k}
\end{equation}

\textbf{Multi-Layer Perceptron (MLP)} with architecture: input layer (55 or 768 features), hidden layer 1 (128 neurons, ReLU, Dropout 0.3), hidden layer 2 (64 neurons, ReLU, Dropout 0.3), output layer (3 neurons, softmax). Trained with Adam optimizer (lr=0.001) and early stopping (patience=20).

Feature importance is analyzed using Random Forest (500 estimators) with incremental validation. ANOVA tests assess between-composer variance for each feature.

\section{Results}\label{sec:results}

\subsection{Classification Performance Comparison}

Table~\ref{tab:performance} shows results across feature configurations.

\begin{table}[h]
\centering
\caption{Classification performance across feature sets. Best result in bold.}
\label{tab:performance}
\begin{tabular}{lrrr}
\toprule
Feature Set & Dimensions & Classifier & Bal. Acc. \\
\midrule
Statistical (12D) & 12 & SVM & 49.3\% \\
Handmade (54D) & 54 & SVM & 63.3\% \\
Handmade (20D, top) & 20 & SVM & \textbf{~67\%} \\
MidiBERT (768D) & 768 & SVM & 47.1\% \\
MidiBERT (768D) & 768 & MLP & 45.3\% \\
\bottomrule
\end{tabular}
\end{table}

The 54D handcrafted feature set significantly outperforms MidiBERT with both classifiers, demonstrating that domain-specific features are more effective than general pretrained representations for this task.

\subsection{Feature Importance Analysis}

Table~\ref{tab:importance} shows the top 10 most important features.

\begin{table}[h]
\centering
\caption{Top 10 most important features for composer classification (54D).}
\label{tab:importance}
\begin{tabular}{rlcl}
\toprule
Rank & Feature & Importance & Category \\
\midrule
1 & f27\_unison\_ratio & 0.0308 & Interval \\
2 & pt\_std & 0.0287 & Texture \\
3 & f3\_pitch\_std & 0.0286 & Pitch \\
4 & f28\_stepwise\_ratio & 0.0280 & Interval \\
5 & mc\_std & 0.0272 & Melody \\
6 & f4\_pitch\_range & 0.0261 & Pitch \\
7 & f22\_staccato\_ratio & 0.0259 & Rhythm \\
8 & f24\_interval\_mean & 0.0243 & Interval \\
9 & f8\_most\_common\_pc\_ratio & 0.0239 & Pitch \\
10 & f25\_interval\_std & 0.0238 & Interval \\
\bottomrule
\end{tabular}
\end{table}

\subsection{Feature Selection Curve}

Figure~\ref{fig:feature_curve} shows the relationship between number of features and classification accuracy. Peak performance (~67\%) is achieved with approximately 20 top-ranked features.

\begin{figure}[h]
\centering
\includegraphics[width=0.9\linewidth]{feature_accuracy_curve_12+54.png}
\caption{SVM feature selection curve showing balanced accuracy as a function of number of top features. Peak performance (~67\%) achieved at ~20 features.}
\label{fig:feature_curve}
\end{figure}

\subsection{Combined Feature Analysis (780D)}

We also evaluated a combined feature set concatenating all features (12D statistical + 54D handmade + 768D MidiBERT). In the combined set, handcrafted features dominate the top rankings, with the first MidiBERT dimension (bert\_dim\_251) appearing only at rank 10. This confirms that for this specific task, domain-specific features encode more relevant information than general pretrained representations.

\subsection{MLP vs. SVM on MidiBERT Embeddings}

We compared SVM and MLP classifiers on the 768D MidiBERT embeddings:

\begin{table}[h]
\centering
\caption{SVM vs. MLP on MidiBERT embeddings.}
\label{tab:mlp}
\begin{tabular}{lc}
\toprule
Classifier & Accuracy \\
\midrule
SVM (RBF) & 47.1\% \\
MLP & 45.3\% \\
\bottomrule
\end{tabular}
\end{table}

SVM outperforms MLP on MidiBERT embeddings, suggesting that the small dataset (264 pieces) is insufficient for MLP to learn effective non-linear boundaries.

\section{Musicological Discussion}\label{sec:discussion}

\subsection{What Features Reveal About Schubert}

Schubert's higher stepwise ratio (f28) aligns with observations about his lyrical, singable melodies that mirror natural speech inflection. Lower texture density reflects his characteristic guitar-like arpeggiated accompaniment (as in ``Die Forelle''). Moderate harmonic complexity captures his expressive but structurally clear chromaticism.

\subsection{What Features Reveal About Schumann}

Schumann's high texture variation (pt\_std) and staccato ratio (f22) reflect his diverse accompaniment patterns and poetic declamation approach. Lower simultaneity suggests arpeggiated textures (as in \textit{Dichterliebe} opening). Complex rhythmic patterns mirror his expressive intensity.

\subsection{What Features Reveal About Brahms}

Brahms' high pitch range (f4) and pitch standard deviation (f3) reflect his symphonic piano writing and rich textures. Conservative interval patterns (f27\_unison\_ratio) support his characterization as classically restrained with folk song influence. Dense chordal writing captures his four-part, string-quartet-like piano voicing.

\subsection{Why Handcrafted Features Outperform MidiBERT}

Domain specificity explains the performance gap: MidiBERT trained on diverse MIDI data may not capture Lieder-specific features. The 264-piece corpus is insufficient for effective fine-tuning of 768-dimensional embeddings. Handcrafted features encode musicological hypotheses about what distinguishes composers.

\subsection{Methodological Rigor: Feature Exclusion}

We explicitly excluded note count and velocity features from our analysis due to concerns about confounding variables and editorial bias. This decision strengthens our methodology by ensuring that classification is based on composer-intentioned patterns rather than artifacts. The fact that our model achieves ~67\% accuracy without these features demonstrates that genuine stylistic markers exist in other musical dimensions (pitch, rhythm, interval, texture).

\section{Exploratory Analysis}\label{sec:exploratory}

\subsection{Case Study: Winterreise}

The 24 songs of Schubert's \textit{Winterreise} show lower variance in texture features compared to his other Lieder, suggesting coherent textural vocabulary throughout the cycle. This quantitatively supports musicological observations about the cycle's unity.

\subsection{Case Study: Dichterliebe}

Schumann's \textit{Dichterliebe} opening shows very low simultaneity and high texture variance, quantitatively confirming the famous arpeggiated introduction. Later songs show increased texture density, reflecting the cycle's emotional progression.

\subsection{Interval Features as Primary Discriminators}

Contrary to our initial hypotheses emphasizing harmonic features, unison ratio (f27) and stepwise ratio (f28) emerged as the most discriminative attributes after note count removal. This suggests that melodic motion patterns carry more stylistic information than deeper tonal properties for same-era composer classification. This finding echoes Youngblood's (1958) original insight about statistical distributions capturing stylistic choice.

\section{Conclusion and Future Work}\label{sec:conclusion}

This exploratory study investigated whether computational features can capture individual composer style in German Lieder. Key findings: (1) Handcrafted features (54D, ~67\%) outperform pretrained embeddings (47.1\% SVM, 45.3\% MLP); (2) Unison ratio, texture variation, and pitch standard deviation are most discriminative; (3) Compact subsets (~20 features) achieve peak performance; (4) Methodological exclusions (note count, velocity) strengthen validity of findings.

Limitations include dataset size (264 pieces), composer coverage (three canonical composers), and feature scope (no chord function analysis). Future work includes chord tonal distance implementation, expanded corpus (Wolf, Strauss, Mahler), multi-modal fusion, and deeper musicologist collaboration.

The goal is not automated attribution but enhanced understanding. By quantifying aspects of musical style, we create tools for asking old questions about what makes each composer distinctive.

\section{AI Usage Statement}

AI tools were used for code documentation and initial draft organization. All experimental design, data analysis, musicological interpretation, and final writing were performed by the authors.

% References
\bibliographystyle{unsrt}
\bibliography{references}

\end{document}
